%% kapa: exam=1 round=31 year=2020 period=2 form=A subject=accounting subject_ko=회계학 seq=15 number=55
%% source: Q-Net 감정평가사 기출문제 (cst00309) artlSeq=5207219 file=extracted/1008/5207219/A형/A형/2020년도 제31회 감정평가사 1차 시험_2교시_A형.pdf
%% conversion: kordoc 4.24.51 -> tools/gen1.py (auto); review=proofread-sonnet; flags=none
\begin{question}{55}
(주)감평의 20x1년 말 예상되는 자산과 부채는 각각 ￦100,000과 ￦80,000으로 부채비율(총부채 $÷$ 주주지분) 400\%가 예상된다. (주)감평은 부채비율을 낮추기 위해 다음 대안들을 검토하고 있다. 다음 설명 중 옳지 않은 것은? (단, (주)감평은 모든 유형자산에 대하여 재평가모형을 적용하고 있다.)
\begin{framedtext}
○ 대안Ⅰ: 토지A 처분(장부금액 ￦30,000, 토지재평가잉여금 ￦1,000, 처분손실 ￦5,000 예상) 후 처분대금으로 차입금 상환 \par
○ 대안Ⅱ: 유상증자(￦25,000) 후 증자금액으로 차입금 상환 \par
○ 대안Ⅲ: 토지B에 대한 재평가 실시(재평가이익 ￦25,000 예상) \par
\end{framedtext}
\begin{choices}
  \item 토지A 처분대금으로 차입금을 상환하더라도 부채비율은 오히려 증가한다.
  \item 토지A를 처분만 하고 차입금을 상환하지 않으면 부채비율은 오히려 증가한다.
  \item 유상증자 대금으로 차입금을 상환하면 부채비율은 감소한다.
  \item 유상증자만 하고 차입금을 상환하지 않더라도 부채비율은 감소한다.
  \item 토지B에 대한 재평가를 실시하면 부채비율은 감소한다.
\end{choices}
\end{question}
\answer{1}
